\chapter{Considerações Finais}

Este trabalho partiu de uma lacuna identificada na interseção entre o
autorregistro de humor e o cotidiano de pessoas com Transtorno Afetivo Bipolar:
embora o \textit{Ecological Momentary Assessment} ofereça um caminho
consolidado para reduzir os vieses de memória e aproximar o registro do momento
em que a experiência ocorre, os aplicativos disponíveis frequentemente
fragmentam sono, humor, gatilhos e ações de cuidado em telas isoladas,
restringem a exportação dos dados e operam como soluções fechadas, com pouca
transparência sobre o uso das informações sensíveis que coletam.

Para responder a essa lacuna, foi desenvolvido o Nexo, uma aplicação móvel de
código aberto que integra esses registros em uma única linha do tempo e produz
indicadores destinados a auxiliar o usuário na percepção de padrões ao longo do
tempo. O sistema é composto por uma API REST e por um aplicativo móvel, ambos
escritos em TypeScript, com processamento de indicadores em segundo plano,
implantação automatizada e uma suíte de testes automatizados concentrada na API,
complementada por testes unitários no aplicativo móvel, conforme descrito ao
longo dos capítulos anteriores.

A pesquisa aplicada de natureza exploratória orientou as decisões de projeto. O
questionário aplicado a 48 potenciais usuários, apresentado no
Capítulo~\ref{chapter:methodologies}, e a análise de aplicativos similares
forneceram os insumos que sustentaram a definição dos requisitos, da identidade
visual e da organização do sistema, bem como a adoção do modelo incremental com
princípios ágeis e \textit{Kanban} no desenvolvimento.

Diante disso, conclui-se que os objetivos propostos foram alcançados: o
protótipo funcional entregue demonstra a viabilidade de uma solução de
autorregistro alinhada aos princípios de privacidade, transparência e autonomia
do usuário sobre os próprios dados.

O trabalho, contudo, apresenta limitações que devem ser explicitadas. A
verificação do aplicativo móvel conta com testes unitários sobre a lógica pura,
mas ainda carece de testes de componente e de integração e não passa pela
integração contínua; assim, regressões na interface podem ser incorporadas ao
código sem que nenhum teste acuse, como discutido no
Capítulo~\ref{cap:testes}. A amostra do questionário é, por conveniência, e não
pode ser generalizada; além disso, o protótipo ainda não foi submetido a uma
avaliação com usuários reais, o que deixa em aberto a confirmação de que as
escolhas de interface e de tom atendem, de fato, ao público-alvo. Por fim,
algumas decisões sensíveis --- como o rótulo de estabilidade apresentado nos
indicadores --- permanecem deliberadamente estáticas, aguardando embasamento
específico antes de se tornarem funcionais.

Os resultados obtidos abrem caminhos de continuidade que, na leitura deste
trabalho, merecem prioridade. O primeiro é a reaplicação da pesquisa junto a
usuários potenciais: repetir o questionário descrito no
Capítulo~\ref{chapter:methodologies} --- ou uma versão adaptada dele --- em uma
nova amostra, desta vez acompanhada de um teste do protótipo funcional,
permitiria comparar as expectativas registradas em 2025 com a percepção
concreta do aplicativo, avaliando usabilidade, tom de comunicação e interesse
pelas funcionalidades, e capturando ajustes antes de uma eventual implantação
mais ampla.

Uma segunda frente é a criptografia no lado do cliente para anotações e
registros de humor, de modo que o conteúdo sensível seja cifrado antes de
deixar o dispositivo e permaneça inacessível mesmo ao operador do servidor.
Essa medida responde diretamente à preocupação com privacidade registrada no
questionário, em que a maioria dos respondentes atribuiu nota máxima à
importância do controle sobre os próprios dados, e reforça o cenário de
auto-hospedagem descrito no Capítulo~\ref{sec:deploy-and-ci}.

A terceira frente é ampliar a testabilidade do aplicativo móvel, cuja suíte
unitária ainda alcança apenas a lógica pura (Seção~\ref{sec:test-mobile}).
Adotar testes de componente e integrar \textit{lint}, verificação de tipos, a
suíte unitária e a suíte \textit{end-to-end} ao fluxo de integração contínua
reduziria o risco de regressões invisíveis e aproximaria o nível de confiança do
\textit{frontend} àquele já alcançado pela API.

Além dessas, as lacunas apontadas na Seção~\ref{sec:test-gaps} --- como a
maquinaria da fila fora do alcance dos testes e a duplicação manual dos dados
de teste --- sugerem ajustes de menor porte, porém de impacto relevante na
manutenibilidade do projeto a médio prazo.

Outra frente que a implementação revelou diz respeito ao ciclo de vida dos
medicamentos. O protótipo trata um regime medicamentoso por um único estado
binário --- em uso ou fora de uso ---, o que se mostrou insuficiente para cobrir
situações corriqueiras do cotidiano do público-alvo: pausar temporariamente um
medicamento por orientação médica, encerrar um tratamento por tempo determinado
e descontinuar definitivamente um medicamento são situações distintas que hoje
não podem ser expressas nem distinguidas pelo sistema. A consequência
alcança também o ciclo de vida do regime: como o sistema guarda apenas o
estado atual, não há como datar quando uma pausa ou um encerramento
ocorreu. O histórico de doses, por outro lado, já não depende disso: o nome
e a dosagem são copiados para cada registro de dose no momento em que ela é
gravada, de modo que editar ou desativar o regime não reescreve
retroativamente o passado. O que falta é representar a vigência do regime
ao longo do tempo, para que o próprio período de uso --- e não apenas a
dose --- seja reconstruível. Representar essa vigência permitiria ainda que a
periodicidade configurada repercutisse de fato na geração dos lembretes, hoje
disparados diariamente qualquer que seja a frequência escolhida.

Outra frente, ainda no campo do que o usuário registra, é a possibilidade de
anexar material aos registros. Hoje o Nexo trabalha apenas com texto e números:
a anotação de um humor é um campo de texto, e o único tipo de mídia que
o sistema guarda é o avatar do perfil. Anexar fotos e gravações de áudio a um
registro de humor --- ou a um gatilho, ou a uma ação de cuidado --- aproximaria
o aplicativo do que o público-alvo já encontra em soluções similares, uma vez
que o Baseline admite entradas por áudio e o DailyBeans permite associar fotos
ao registro. A mudança, porém, não é apenas de interface. Parte da
infraestrutura já existe e funciona: o envio do avatar exercita o caminho
completo de um arquivo, do formulário multipart no aplicativo até
o armazenamento de objetos, com a chave e a URL gravadas no registro
correspondente. O que falta é desenhar um fluxo para um conteúdo que não
é único nem substituível como o avatar: um anexo se acumula, tem vínculo com um
registro específico e precisa sobreviver à edição dele. Isso exigiria um modelo
de dados para o anexo, com tipo, tamanho, duração e o vínculo de origem; um
envio que aceite mais de um arquivo e outras mídias além de imagem, hoje
barradas pelo filtro do \textit{endpoint} de avatar; e a inclusão dos anexos na
exportação integral, que se pretende um retrato completo dos dados do usuário.
Acrescenta ainda uma questão de custo e de retenção: áudio e imagem ocupam
muito mais espaço que texto, e o armazenamento de objetos do projeto opera hoje
com um único \textit{bucket} de leitura pública, adequado a avatares, mas não
a conteúdo íntimo.

Dos anexos de áudio decorre uma segunda possibilidade: transcrevê-los no
próprio dispositivo, e não no servidor. Modelos abertos de porte reduzido,
executados por \textit{runtimes} que cabem em um aplicativo móvel, já tornam
viável transcrever alguns minutos de fala sem que o áudio saia do aparelho.
O uso mais imediato seria o de uma anotação falada: o usuário gravaria o relato
e pediria a transcrição, que passaria a preencher o campo de texto daquele
humor ou daquele gatilho. A motivação não é apenas de conveniência. Manter
o som no dispositivo dialoga com a preferência, registrada no questionário, de
guardar os dados localmente, e sobretudo com a frente de criptografia no
cliente discutida anteriormente --- uma anotação cifrada no aparelho não pode
ser transcrita no servidor, de modo que as duas frentes se reforçam em vez de
competir. Em contrapartida, o tamanho do modelo a baixar, a latência
e o consumo de bateria em aparelhos intermediários, a qualidade da transcrição
em português e o próprio escoamento do texto produzido, já que o aplicativo
ainda não possui fila de envio \textit{offline}, são obstáculos que só uma
avaliação prática resolveria. Registre-se que a transcrição é uma
funcionalidade autônoma, útil mesmo que nenhuma outra frente avance: ela serve
tanto ao registro avulso quanto, mais adiante, ao fornecimento de texto a um
fluxo conduzido por perguntas.

A forma de preencher o registro é outra frente que a implementação sugeriu,
e diz respeito sobretudo à interface. O fluxo atual é declaradamente um
formulário: escolhe-se o humor, julgam-se os níveis de energia, estresse
e ansiedade em uma escala de 1 a 10, selecionam-se sentimentos e intensidades
e, por fim, escreve-se uma anotação, em uma sequência fixa de telas. Uma
alternativa seria conduzir o registro como uma conversa, perguntando como foi
o dia e quanto de estresse ele exigiu, e adaptando a próxima pergunta ao que
foi respondido --- um fluxo guiado, em vez de um conjunto de campos. Vale notar
que os dados já existentes no modelo suportariam essa mudança sem alteração de
esquema, o que torna a frente uma questão de apresentação e de roteiro, e não
de domínio. O cuidado é que um fluxo conversacional pode ficar mais lento que
o formulário que pretende substituir: o registro por toques resolve-se em
poucos segundos e é justamente o que boa parte do público procura. A condução
por perguntas corre ainda o risco de soar avaliativa --- algo que este trabalho
evitou deliberadamente ao não emitir juízos sobre o estado do usuário ---, de
modo que a redação das perguntas pesaria tanto quanto a mecânica da tela. Uma
saída plausível seria oferecer os dois caminhos, deixando ao usuário a escolha
entre o registro rápido e o conduzido, conforme o momento.

Por fim, os relatórios comportam ao menos duas ampliações. A primeira é de
conteúdo: o documento desta versão cobre apenas os dados de humor, e a intenção
de incorporar o sono, as atividades e os padrões de gatilho, incluindo
a correlação entre sono e energia já calculada nos indicadores, foi adiada.
A segunda é de forma: hoje o usuário escolhe apenas o período, e todo
o restante do documento é fixo. Permitir escolher o que entra no relatório, ou
oferecer modelos prontos para finalidades distintas --- um resumo curto para
uma consulta, um panorama mais longo para revisão pessoal ---, tornaria
o artefato mais expressivo e mais próximo do uso real, que é levar algo legível
a um profissional de saúde. O risco é que a ampliação agrave o problema de
interpretação já registrado para os indicadores: um gráfico ou um recorte novo
precisa declarar sobre qual janela foi calculado e sobre qual escala repousa,
sob pena de sugerir conclusões que o dado não sustenta --- o mesmo cuidado que
levou a manter estático o rótulo de estabilidade. Uma versão configurável,
portanto, exigiria que cada elemento do relatório carregasse sua própria
legenda de leitura.
